\chapter{Modélisation UML}

Ce chapitre présente les diagrammes UML de la solution Droppy. Les fichiers sources PlantUML se trouvent dans le dossier \texttt{diagrams/}. Pour générer les images, exécutez le script \texttt{generate-diagrams.bat} ou utilisez \url{https://www.plantuml.com/plantuml}.

\section{Diagramme de cas d'utilisation}

\textbf{Acteurs identifiés :}
\begin{itemize}
  \item \textbf{Opérateur} --- consulte le dashboard, valide les alertes ;
  \item \textbf{Planificateur (Système)} --- déclenche l'analyse périodique ;
  \item \textbf{Microservice IA (FastAPI)} --- acteur secondaire (système externe).
\end{itemize}

\umlfigure{usecase}{Diagramme de cas d'utilisation — Plateforme Droppy}

\section{Diagramme de classes}

Le diagramme de classes modélise les entités Java (backend), les services Angular (frontend) et les modules Python (microservice IA).

\umlfigure{class_diagram}{Diagramme de classes — Architecture complète Droppy}

\section{Environnement de travail}

La réalisation de la solution \droppy{} s'est appuyée sur un environnement
technologique hétérogène, dans lequel chaque composant repose sur des
technologies et des outils adaptés à son rôle.

\subsection{Environnement matériel}

Le développement et les tests ont été réalisés sur un poste de travail
présentant la configuration suivante.

\begin{table}[H]
\centering
\caption{Configuration matérielle du poste de développement}
\label{tab:environnement-materiel}

\renewcommand{\arraystretch}{1.35}

\begin{tabularx}{\textwidth}{
    >{\bfseries}p{3.6cm}
    >{\raggedright\arraybackslash}X
}

\toprule
Élément & Caractéristique \\
\midrule

Processeur &
Intel Core i7 (ou équivalent), 11e génération ou supérieure. \\
\\[2pt]

Mémoire vive &
16 Go de RAM. \\
\\[2pt]

Stockage &
Disque SSD de 512 Go. \\
\\[2pt]

Système d'exploitation &
Windows 10/11, utilisé pour le développement local. \\
\\[2pt]

Réseau &
Accès réseau local pour les échanges entre les services (FastAPI, Spring
Boot, Angular) et les équipements IoT. \\

\bottomrule

\end{tabularx}
\end{table}

\subsection{Environnement logiciel}

L'environnement logiciel regroupe les outils de développement, le
versionnement du code ainsi que les langages et frameworks utilisés par
chaque composant de la solution.

\subsubsection{Frameworks et technologies}

Le tableau~\ref{tab:environnement-travail} récapitule les principaux
frameworks et technologies mobilisés pour le développement de la solution.

\begin{table}[H]
\centering
\caption{Frameworks et technologies de la solution Droppy}
\label{tab:environnement-travail}

\renewcommand{\arraystretch}{1.35}

\begin{tabularx}{\textwidth}{
    >{\bfseries}p{3.6cm}
    >{\raggedright\arraybackslash}X
}

\toprule
Composant & Technologies et outils \\
\midrule

Système d'exploitation &
Développement et exécution locale sous Windows. \\
\\[2pt]

Frontend (\texttt{dashboard-angular}) &
Angular~19, TypeScript~5.7, RxJS~7.8 et Node.js (version LTS). \\
\\[2pt]

Backend (\texttt{core-service}) &
Java~17, Spring Boot~3.3.1, Maven et Liquibase pour la gestion des
migrations de schéma. \\
\\[2pt]

Microservice IA (\texttt{src}) &
Python~3.10+, FastAPI, scikit-learn (Isolation Forest), joblib
(sérialisation des modèles), pandas, numpy, SQLAlchemy et uvicorn. \\
\\[2pt]

Base de données &
PostgreSQL en environnement réel, H2 pour les tests. \\
\\[2pt]

Modèles &
Isolation Forest et modèle de prévision quotidienne (régression Ridge),
sauvegardés au format joblib dans \texttt{models/}. \\
\\[2pt]

Communication &
REST/HTTP entre les services et WebSocket STOMP pour les notifications
temps réel vers le dashboard. \\
\\[2pt]

Outils de développement &
Git/GitHub pour le versionnement, IDE (IntelliJ IDEA, Visual Studio Code),
PlantUML pour la modélisation des diagrammes. \\
\\[2pt]

Tests et validation &
pytest pour le pipeline Python, JUnit pour le backend Spring, frameworks
de test Angular (Jasmine/Karma) pour le frontend. \\

\bottomrule

\end{tabularx}
\end{table}

Cette organisation permet de développer chaque composant de manière
indépendante tout en maintenant des interfaces claires entre les services.

\section{Diagrammes de séquence}

Les diagrammes de séquence décrivent les échanges entre les différents
composants de la solution pour les principaux scénarios d'utilisation.

\subsection{Détection automatique planifiée}

Ce scénario illustre le déclenchement périodique de l'analyse par le
planificateur. Le microservice IA applique l'ensemble des mécanismes de
détection, puis le backend persiste les nouvelles anomalies et notifie le
dashboard en temps réel via WebSocket.

\umlfigure{sequence_scheduled}{Séquence — Détection automatique planifiée d'anomalies}

\subsection{Validation d'une anomalie}

Ce scénario décrit la validation d'une anomalie par l'opérateur depuis le
dashboard. Le backend met à jour le statut de l'anomalie en base et diffuse
la modification aux clients connectés, le tableau étant ensuite rafraîchi.

\umlfigure{sequence_validate}{Séquence — Validation d'une anomalie par l'opérateur}

\subsection{Entraînement des modèles ML}

Ce scénario illustre l'entraînement des modèles associés à un dispositif.
Les données historiques sont chargées, nettoyées et prétraitées, puis
l'Isolation Forest et le modèle de prévision sont entraînés et sérialisés
au format joblib.

\umlfigure{sequence_train}{Séquence — Entraînement des modèles ML}

\subsection{Consultation filtrée des anomalies}

Ce scénario décrit la consultation des anomalies par l'opérateur avec
application de filtres (MAC, statut, sévérité). La requête paginée est
exécutée par le backend, puis les résultats sont affichés dans le dashboard.

\umlfigure{sequence_get}{Séquence — Consultation filtrée des anomalies}

\section{Code PlantUML}

Les sources complètes des diagrammes sont disponibles en annexe (Listing~\ref{lst:plantuml-usecase}) et dans le dossier \texttt{diagrams/}.

\lstinputlisting[caption={Diagramme de cas d'utilisation (PlantUML)}, label={lst:plantuml-usecase}, language=]{diagrams/usecase.puml}